\ficha{Fora do plenário em 1906: governo, Murtinho e a baixa de setembro}

\bloco{Fontes lidas}

\textit{Correio da Manhã} de 21 e 29 de setembro e de 10 e 16 de outubro de
1906; \textit{Gazeta de Notícias} de 24 e 26 de setembro e de 23 de outubro de
1906; discurso de David Campista de 5 de outubro (\textit{Documentos
Parlamentares}, p.~538). Como retrospecto, o discurso de Pinheiro Machado no
Senado em dezembro de 1913, no resumo de \textit{O Paiz} e na transcrição do
\textit{Correio da Manhã}, ambos de 28 de dezembro de 1913.

\bloco{Síntese}

Os anais registram o que se disse na tribuna. A imprensa de setembro e outubro
de 1906 registra três fatos que os anais só deixam entrever: o governo
Rodrigues Alves era contrário ao projeto; a queda do câmbio de 19 de setembro
foi atribuída à notícia de uma emenda de Alcindo Guanabara; e Joaquim Murtinho
renunciou ao mandato de senador por causa da Caixa. Sete anos depois, Pinheiro
Machado conta a mesma história como memória política da criação.

\bloco{O governo contra o projeto}

O próprio relator o admite na tribuna, referindo-se ao presidente e ao ministro
da Fazenda: \tr[P:1:544]{Sei bem que SS. EEx. não estão de accôrdo com a
reforma} \dpa{p.~538}. O \textit{Correio da Manhã} passa em revista os
ex-ministros da Fazenda da República e informa que Rodrigues Alves e Leopoldo de
Bulhões

\cit[I:per089842_1906_01916:1]{declararam a sua franca opposição á politica das
aventuras financeiras iniciada pelo Convenio de Taubaté}{\jor{cm19061010}{Correio
da Manhã}{10 out. 1906}{1}}

e conclui que, dos dez, \tr[I:per089842_1906_01916:1]{o unico que está
inequivocamente favoravel ao projecto da Caixa de Conversão, é o sr. Bernardino
de Campos}. A \textit{Gazeta de Notícias} explica a reação de Londres pelo fato
de os capitalistas estarem \tr[I:per103730_1906_00267:7]{ignorando como se acham
divorciados nesta materia o poder legislativo e o executivo}
(\jor{gn19060924}{Gazeta de Notícias}{24 set. 1906}{7}).

\bloco{A baixa de 19 de setembro e a emenda de Alcindo}

Em 25 de setembro Alcindo Guanabara menciona na tribuna a imputação de ter
provocado a baixa do câmbio (ficha 1). A imprensa confirma que a imputação
existiu. O \textit{Correio da Manhã}, em texto de primeira página, defende o
Banco do Brasil e o governo e inverte a culpa. A votação do projeto e as
emendas anunciadas por Alcindo teriam mostrado que no Congresso
\tr[I:per089842_1906_01897:1]{soprava um sinistro vento de vesania}:

\cit[I:per089842_1906_01897:1]{Quer-se fazer um salto ainda maior para o abysmo
com a taxa de 12}{\jor{cm19060921}{Correio da Manhã}{21 set. 1906}{1}}

O mesmo texto chama o plano de Taubaté de \tr[I:per089842_1906_01897:1]{aventura
sem precedentes}. A resenha semanal do mercado de câmbio da \textit{Gazeta de
Notícias} lista as versões correntes para a queda e adota a seguinte: informados
de que Alcindo apresentaria emendas quebrando o padrão a 12, os capitalistas de
Londres

\cit[I:per103730_1906_00267:7]{temeram a possibilidade de se tornar lei a emenda
Alcindo}{\jor{gn19060924}{Gazeta de Notícias}{24 set. 1906}{7}}

Dois dias depois o jornal noticia que o deputado \tr[I:per103730_1906_00269:1]{resolveu
não apresentar o seu trabalho} (\jor{gn19060926}{Gazeta de Notícias}{26 set.
1906}{1}), o que coincide com o discurso de 25 de setembro, em que Alcindo diz
não ter apresentado emenda.

\bloco{A renúncia de Joaquim Murtinho}

O \textit{Correio da Manhã} anuncia a renúncia no fim de setembro e lhe dá o
motivo:

\cit[I:per089842_1906_01905:1]{por ser o sr. Murtinho contrario á approvação do
projecto que crea a Caixa de Conversão}{\jor{cm19060929}{Correio da Manhã}{29
set. 1906}{1}}

Em 16 de outubro o jornal descreve a divisão entre os chefes do Bloco: de um
lado Murtinho e Rui Barbosa, do outro Afonso Pena e Jorge Tibiriçá, e
acrescenta: \tr[I:per089842_1906_01922:1]{O sr. Pinheiro Machado hesita}
(\jor{cm19061016}{Correio da Manhã}{16 out. 1906}{1}). A \textit{Gazeta de
Notícias} noticia a renúncia em 22 de outubro, e a edição do dia seguinte traz a
moção de pesar aprovada pelo Senado. O jornal comenta, em texto de primeira
página:

\cit[I:per103730_1906_00296:1]{o Sr. Murtinho combate a Caixa de Conversão, por
omissão}{\jor{gn19061023}{Gazeta de Notícias}{23 out. 1906}{1}}

\bloco{Rui Barbosa antes do voto}

Rui vota a favor em 26 de novembro (ficha 2). Em outubro, o \textit{Correio da
Manhã} lhe atribuía posição diversa: \tr[I:per089842_1906_01916:1]{são notorias
as suas radicaes divergencias com o projecto Campista}
(\jor{cm19061010}{Correio da Manhã}{10 out. 1906}{1}). A frase é do jornal, que
escrevia contra o projeto. Não é fala de Rui.

\bloco{A memória de Pinheiro Machado, 1913}

Em dezembro de 1913, no Senado, Pinheiro Machado reconta a criação da Caixa.
Diz que o projeto teve contra si

\cit[I:per178691_1913_10674:1]{não só a acção directa do governo, como ainda de
orgãos importantes da imprensa desta capital}{\jor{opaiz19131228}{O Paiz}{28
dez. 1913}{1}}

e nomeia os jornais: \tr[I:per178691_1913_10674:1]{o Jornal do Commercio, o
Paiz, o Correio da Manhã e outros}. Conferi a lista na imagem de \textit{O
Paiz}. No \textit{Correio da Manhã} do mesmo dia, que transcreve o discurso em
primeira pessoa, o trecho tem um espaço em branco, com pontos residuais, entre
\textit{Jornal do Commercio} e \textit{o Correio da Manhã}, também conferido na
imagem; não sei explicar a lacuna. Pinheiro diz ainda que Murtinho renunciou
\tr[I:per089842_1913_05446:3]{por ser infenso á Caixa de Conversão}
(\jor{cm19131228}{Correio da Manhã}{28 dez. 1913}{3}) e lê a resposta que lhe
deu em 1906:

\cit[I:per178691_1913_10674:2]{A estabilidade do cambio não demora, ao
contrario, apressa a conversão}{\jor{opaiz19131228}{O Paiz}{28 dez. 1913}{2}}

\bloco{Achados}

\begin{enumerate}[leftmargin=*,nosep]
\item A renúncia de Murtinho por causa da Caixa é fato de outubro de 1906,
documentado no mesmo mês por dois jornais. A versão de Pinheiro Machado em 1913
se confirma nesse ponto, e também quanto à sua própria hesitação.
\item A oposição do Executivo ao projeto está dita por Campista na tribuna e
pela imprensa. A Caixa foi criada pelo Congresso contra a posição declarada do governo que
terminava.
\item A imputação a Alcindo é datada: queda do câmbio em 19 de setembro,
acusação em 21 de setembro, desistência do substitutivo noticiada em 26 de
setembro.
\item A posição atribuída a Rui Barbosa antes do voto vem só de jornal
adversário do projeto.
\end{enumerate}

\bloco{Limites}

Não localizei o discurso de Bulhões em banquete nem a mensagem de Rodrigues
Alves que o \textit{Correio da Manhã} menciona. A memória de Pinheiro Machado é
retrospectiva e interessada: ele fala como chefe político que reivindica a
autoria da reforma. O \textit{Jornal do Commercio} não integra o corpus.
